% !TEX program = pdflatex
\documentclass[a4paper,12pt]{article}
\usepackage[T1]{fontenc}
\usepackage[utf8]{inputenc}
\usepackage[czech]{babel}
\usepackage{lmodern}
\usepackage{tul-semestralni}

% Změň pouze kód: FS, FT, FP, EF, FUA, FM nebo FZS.
\Fakulta{FS}
\newcommand{\Predmet}{Modelování a Simulace}
\newcommand{\NazevPrace}{Semestrální práce 3}
\newcommand{\AutorPrace}{Adam Kábrt}
\newcommand{\DatumPrace}{\today}
\newcommand{\RokPrace}{\the\year}
\newcommand{\VerzeZadani}{2}
\newcommand{\PouzitySW}{Creo Simulate}

\begin{document}
\TitulniStrana
% Titulní strana se nepočítá; text začíná stránkou 1.
\section{Zadání a cíl práce}
Zde stručně popište zadání, řešený problém a cíl semestrální práce.

\section{Postup řešení}
Uveďte použité metody, vstupní údaje a předpoklady výpočtu.
Na použitý zdroj odkažte citací, například \cite{ukazka}.
\subsection{Výpočet}
Rovnice lze zapisovat například takto:
\begin{equation}
  \sigma = \frac{F}{A},
\end{equation}
kde $F$ je zatěžující síla a $A$ plocha průřezu.

\section{Výsledky}
Výsledky doplňte obrázky, grafy nebo tabulkami. Obrázky ukládejte
do složky projektu a vkládejte pomocí příkazu \verb|\includegraphics|.
\begin{table}[htbp]
  \centering
  \caption{Ukázka přehledu výsledků}
  \begin{tabular}{lrr}
    \hline
    Veličina & Hodnota & Jednotka \\
    \hline
    Síla & 1000 & N \\
    Plocha & 100 & mm$^2$ \\
    Napětí & 10 & MPa \\
    \hline
  \end{tabular}
\end{table}

\newpage
\section{Závěr}
Shrňte hlavní výsledky a zhodnoťte splnění cíle práce.

% Zdroje přidávejte do literatura.bib ve stejné složce jako tento soubor.
% V textu použijte \cite{klic}; seznam se vytvoří automaticky.
\bibliographystyle{unsrt}
\bibliography{literatura}
\end{document}
